% ch2_e 第2章 §2.3尾–§2.4 (书页 58–67 / p073–p082)
%--D-TAIL--
% 本块顶部为问题 2.4.2 的第 6 问（题干与小问 1–5 在 ch2_d 末尾），装配器会接回第 5 问之后。
\begin{enumerate}
\item[6.] 讨论对 $g$ 与 $m_0$ 的哪些取值，半经典隧穿（tunneling）图像才是对亚稳态（meta-stable state）衰变的一个好的描述。
\end{enumerate}
%--D-TAIL 结束--

\subsection{摩擦的量子理论}

\begin{keypoints}
\item 量子系统中的耗散（摩擦）可以通过让系统与一个热库（heat bath，即一组谐振子）耦合来模拟。
\end{keypoints}

粒子受到的摩擦来自它与环境的耦合。由于这种耦合，粒子会把能量丢失到环境中。因此，要建立摩擦的量子理论，我们必须把与环境的耦合包含进来。这里我们用一组谐振子来模拟环境（Caldeira and Leggett, 1981）。我们需要搭建一个能给出有限摩擦系数（friction coefficient）的恰当环境。这使我们能够理解存在摩擦时粒子的量子运动。

我们的模型由如下路径积分描述：
\begin{equation*}
Z = \text{常数} \times \int \mathcal{D}[x(t)]\,\mathcal{D}[h_i(t)]\;\mathrm{e}^{\mathrm{i}\int \mathrm{d}t\,\left(\frac{1}{2}m\dot{x}^2 + \sum_i \frac{1}{2}\dot{h}_i^2 - \frac{1}{2}(\Omega_i h_i + g_i x)^2\right)}
\end{equation*}
（译注：原书此式指数中第一项误排为 $\frac{1}{2}m\dot{x}$，遗漏了平方；现按式 (2.4.8) 改正为 $\frac{1}{2}m\dot{x}^2$。）

其中 $x$ 是粒子的坐标，诸 $h_i$ 描述一簇谐振子。注意我们的系统具有平移对称性 $x \to x + \text{常数}$。把 $h_i$ 积掉之后，我们得到（见 2.2.3 节）
\begin{equation*}
Z = \text{常数} \times \int \mathcal{D}[x(t)]\;\mathrm{e}^{\mathrm{i}\int \mathrm{d}t\,\frac{m}{2}\dot{x}^2 + \mathrm{i}\int \mathrm{d}t\,\mathrm{d}t'\,\frac{1}{4}G_{\mathrm{os}}(t-t')(x(t)-x(t'))^2} \tag{2.4.8}
\end{equation*}
其中
\begin{align*}
G_{\mathrm{os}}(t-t') &= \sum_i \Omega_i^2 g_i^2 (-\mathrm{i})\langle h_i(t)h_i(0)\rangle = -\mathrm{i}\sum_i \frac{g_i^2\Omega_i}{2}\,\mathrm{e}^{-\mathrm{i}|t-t'|\Omega_i}\\
&= -\mathrm{i}\int_0^{\infty} \mathrm{d}\Omega\; n(\Omega)\,\frac{\Omega g^2(\Omega)}{2}\,\mathrm{e}^{-\mathrm{i}|t-t'|\Omega},
\end{align*}
$n(\Omega) \equiv \sum_i \delta(\Omega - \Omega_i)$ 是态密度（density of states），而 $g^2(\Omega) \equiv \sum_i g_i^2\delta(\Omega - \Omega_i)/\sum_i \delta(\Omega - \Omega_i)$。式 (2.4.8) 中的双时作用量（double-time action）给出如下运动方程：
\begin{equation*}
-\omega^2 m x_\omega = -(G_{\mathrm{os}}(\omega) - G_{\mathrm{os}}(0))\, x_\omega,
\end{equation*}
它可以改写为\footnote{这里我们用了 $\mathrm{Im}G_{\mathrm{os}}(\omega = 0) = 0$。注意 $\mathrm{Im}D_{\mathrm{os}}(\omega) = -\mathrm{Im}D_{\mathrm{os}}(-\omega)$，因而 $\mathrm{Im}D_{\mathrm{os}}(\omega = 0) = 0$。由于 $\mathrm{Im}G_{\mathrm{os}}(\omega) = \mathrm{sgn}(\omega)\mathrm{Im}D_{\mathrm{os}}(\omega)$，我们发现 $\mathrm{Im}G_{\mathrm{os}}(\omega = 0) = 0$。}
\begin{align*}
-\omega^2 m^* x_\omega &= -(-\mathrm{i}\omega)\gamma x_\omega, \tag{2.4.9}\\
m^* &= m - \frac{\mathrm{Re}(G_{\mathrm{os}}(\omega) - G_{\mathrm{os}}(0))}{\omega^2}, \qquad \gamma = -\frac{\mathrm{Im}G_{\mathrm{os}}(\omega)}{\omega}
\end{align*}
若 $m^*$ 与 $\gamma$ 在小 $\omega$ 下有限，则式 (2.4.9) 变为 $m^*\frac{\mathrm{d}^2x}{\mathrm{d}t^2} = -\gamma\frac{\mathrm{d}x}{\mathrm{d}t}$，它描写一个质量为 $m^*$ 的粒子，其摩擦由摩擦系数 $\gamma$ 描述。我们看到，与一簇谐振子的耦合可以在粒子上产生摩擦。摩擦（因而是耗散）来自关联的虚部 $\mathrm{Im}G_{\mathrm{os}}(\omega)$。这种耦合还把粒子的质量从 $m$ 改变为 $m^*$。这里 $m^*$ 称为粒子的有效质量（effective mass）。修正量 $m^* - m$ 来自关联的实部 $\mathrm{Re}G_{\mathrm{os}}(\omega)$。物理上，粒子使谐振子发生形变，形变随粒子一起运动，从而改变了粒子的有效质量。

上面的摩擦量子理论有一个问题。由于 $\mathrm{Im}G_{\mathrm{os}}(\omega) < 0$，我们看到摩擦系数 $\gamma(\omega) > 0$（当 $\omega > 0$），这是正确的；然而当 $\omega < 0$ 时 $\gamma(\omega) < 0$，这似乎毫无道理。我们注意到，式 (2.4.8) 中的有效作用量在时间反演（time reversal）$t \to -t$ 下不变。其后果是：如果 $\omega > 0$ 的解对应于衰减过程，那么 $\omega < 0$ 的解就对应于衰减过程的逆过程。

要真正地描述摩擦，我们需要一个不具有时间反演不变性的形式框架。为此，重要的是注意到：在存在摩擦时，$t = -\infty$ 处的态（即 $|0_-\rangle$）与 $t = \infty$ 处的态（即 $|0_+\rangle$）是非常不同的，因为随着粒子慢下来，谐振子被激发到越来越高的态。这使得我们的路径积分 $Z = \langle 0|U(\infty,-\infty)|0\rangle$ 成为一个糟糕的出发点，因为它等于零。于是，这里的问题是为非平衡系统构造路径积分。闭合时间路径积分（closed-time path integral），或称 Schwinger–Keldysh 形式（Schwinger, 1961; Keldysh, 1965; Rammer and Smith, 1986），是解决这些问题的一条途径。Schwinger–Keldysh 形式从一个不同的路径积分 $Z_{\mathrm{close}} = \langle 0|U(-\infty,\infty)U(\infty,-\infty)|0\rangle$ 出发，它从 $t = -\infty$ 走到 $t = +\infty$ 再返回 $t = -\infty$。

上面发展的双时路径积分（double-time path integral）(2.4.8) 只能描写耗散系统的平衡性质，例如存在摩擦/耗散时围绕基态的涨落。它不能描写非平衡性质，例如由摩擦引起的速度减慢，因为这涉及 $t = \pm\infty$ 处两个不同的态 $|0_-\rangle$ 与 $|0_+\rangle$。

这里遇到的问题与 2.2 节中遇到的问题类似，在那里我们曾试图用路径积分计算响应率的虚部——它同样与耗散相联系。在那里，问题是靠把编时关联函数换成响应函数而解决的。这里的问题可以用同样的方式解决（或者更确切地说，用手工方式加以修正）。如果在\emph{运动方程}中把 $G_{\mathrm{os}}(\omega)$ 换成 $D_{\mathrm{os}}(\omega)$（注意 $D_{\mathrm{os}}(t) = 2\Theta(t)\mathrm{Re}G_{\mathrm{os}}(t)$）——但在双时路径积分中不作这种替换——那么摩擦系数 $\gamma$ 将由
\begin{equation*}
\gamma = -\frac{\mathrm{Im}D_{\mathrm{os}}(\omega)}{\omega} \tag{2.4.10}
\end{equation*}
给出，它对 $\omega > 0$ 与 $\omega < 0$ 都总是正的。在记住上述限制的前提下，我们可以说双时路径积分 (2.4.8) 能够描述一个有摩擦的系统。摩擦系数由式 (2.4.10) 给出。

由于
\begin{equation*}
G_{\mathrm{os}}(\omega) = \int_0^{\infty} \mathrm{d}\Omega\;\frac{\Omega^2 n(\Omega)g^2(\Omega)}{\omega^2 - \Omega^2 + \mathrm{i}0^+},
\end{equation*}
我们求得 $\mathrm{Im}G_{\mathrm{os}}(\omega) = -\frac{\pi}{2}|\omega|n(|\omega|)g^2(|\omega|)$。利用 $\mathrm{Im}D_{\mathrm{os}}(\omega) = \mathrm{sgn}(\omega)\mathrm{Im}G_{\mathrm{os}}(\omega)$ 这一事实，我们求得对小 $\omega$ 值
\begin{equation*}
\gamma = \frac{\pi}{2}\, n(|\omega|)\, g^2(|\omega|) \tag{2.4.11}
\end{equation*}
因此，要使 $\gamma$ 在小 $\omega$ 下有限，我们要求当 $\Omega \to 0$ 时 $n(\Omega)g^2(\Omega)$ 有限。

让我们假设当 $\Omega < \Omega_0$ 时 $n(\Omega)g^2(\Omega) = n_0 g_0^2$，而当 $\Omega > \Omega_0$ 时 $n(\Omega)g^2(\Omega) = 0$。我们求得
\begin{align*}
G_{\mathrm{os}}(\omega) &= -n_0 g_0^2 \Omega_0 + \frac{n_0 g_0^2|\omega|}{2}\ln\left|\frac{\Omega_0 + |\omega|}{\Omega_0 - |\omega|}\right| - \mathrm{i}\,\frac{\pi}{2}|\omega|\, n_0 g_0^2\\
G_{\mathrm{os}}(t) &= \mathrm{i}\,\frac{n_0 g_0^2}{2t^2}\left(1 - (\mathrm{i}|t|\Omega_0 + 1)\,\mathrm{e}^{-\mathrm{i}|t|\Omega_0}\right)
\end{align*}
除了有限的摩擦系数 $\gamma = \frac{\pi}{2}n_0 g_0^2$ 之外，我们还得到有效质量 $m^* = m - \frac{n_0 g_0^2}{\Omega_0^2}$。

我们想指出，$G_{\mathrm{os}}(t)$ 中的 $(\mathrm{i}|t|\Omega_0 + 1)\mathrm{e}^{-\mathrm{i}|t|\Omega_0}$ 项来自 $n(\Omega)g^2(\Omega)$ 在 $\Omega = \Omega_0$ 处的不连续性。如果 $n(\Omega)g^2(\Omega)$ 是 $\Omega$ 的光滑函数，那么这一项在大 $t$ 极限下将指数式地消失。例如，若选取 $n(\Omega)g^2(\Omega) = n_0 g_0^2\,\mathrm{e}^{-\Omega^2/\Omega_0^2}$，则在大 $t$ 极限下将有 $G_{\mathrm{os}}(t) = \mathrm{i}\frac{n_0 g_0^2}{2t^2}$。摩擦系数仍由 $\gamma = \mathrm{sgn}(\omega)\frac{\pi}{2}n_0 g_0^2$ 给出。采用这样的 $n(\Omega)g^2(\Omega)$ 选择之后，我们发现一个具有摩擦 $\gamma = \frac{\pi}{2}n_0 g_0^2$ 的粒子由如下双时作用量描述：
\begin{equation*}
S = \int \mathrm{d}t\;\frac{m}{2}\dot{x}^2 + \mathrm{i}\int \mathrm{d}t\,\mathrm{d}t'\;\frac{\gamma\,(x(t) - x(t'))^2}{4\pi(t - t')^2} \tag{2.4.12}
\end{equation*}

%FIG{2.10}: {一个 RCL 电路。}

\subsection*{问题 2.4.3}

\textbf{有效质量：}
\begin{enumerate}
\item[1.] 假设对小 $\Omega$ 有 $n(\Omega)g^2(\Omega) \propto \Omega^\alpha$，且当 $\Omega > \Omega_0$ 时 $n(\Omega)g^2(\Omega) = 0$。求出使有效质量 $m^*$ 发散的 $\alpha$ 值。
\item[2.] 考虑 $n(\Omega)g^2(\Omega) = N_0\,\mathrm{e}^{-\Omega/\Omega_0}$ 与 $n(\Omega)g^2(\Omega) = N_0\,\mathrm{e}^{-(\Omega/\Omega_0)^2}$ 两种情形。哪一种情形的有效质量发散？
\end{enumerate}

\subsection{RCL 电路的量子理论}

\begin{keypoints}
\item RCL 电路的量子动力学与一个带摩擦粒子的量子动力学类似。
\item 如何在一个量子 RCL 电路中纳入电荷量子化（charge quantization）。
\end{keypoints}

RCL 电路是一个简单的电路，我们知道如何在经典层次上描述它。既然一切都被认为是用量子理论描述的，问题就产生了：如何用量子力学来描述一个 RCL 电路？

让我们先考虑电容为 $C$ 的一个理想电容器。电容器的态由其电荷 $q$ 描述；如果我们忽略电荷量子化，$q$ 是一个实数。理想电容器的量子理论很简单：希尔伯特空间由态 $\{|q\rangle\}$ 张成，哈密顿量为 $H = \frac{\hat{q}^2}{2C}$，其中电荷算符 $\hat{q}$ 定义为 $\hat{q}|q\rangle = q|q\rangle$。于是，一个量子电容器等价于一条直线上的自由粒子，$q$ 是它的“动量”。哈密顿量 $H = \frac{\hat{q}^2}{2C}$ 的形式告诉我们，这个粒子的质量恰为 $C$。相应的经典拉格朗日量为 $L = \frac{1}{2}C\dot{x}^2$。

为了描述一个 CL 电路，我们加入一个势能项，考虑
\begin{equation*}
L_{\mathrm{CL}} = \frac{C}{2}\dot{x}^2 - \frac{1}{2L}x^2 \tag{2.4.13}
\end{equation*}
由此得到的运动方程为
\begin{equation*}
C\frac{\mathrm{d}^2x}{\mathrm{d}t^2} = -\frac{x}{L}
\end{equation*}
如果我们把 $x/L$ 解释为电流 $I$、把 $L$ 解释为电感，它与 CL 电路的方程
\begin{equation*}
V = \frac{q}{C} = L\frac{\mathrm{d}I}{\mathrm{d}t}, \qquad\qquad I = -\frac{\mathrm{d}q}{\mathrm{d}t},
\end{equation*}
等价。我们看到，CL 电路等价于一个坐标为 $x$、质量为 $C$、弹簧常数为 $1/L$ 的谐振子。于是，量子 CL 电路由哈密顿量
\begin{equation*}
H = \frac{1}{2C}\hat{q}^2 + \frac{1}{2L}\hat{x}^2 \tag{2.4.14}
\end{equation*}
描述，其中 $\hat{q}$ 是电荷算符；它也是与 $\hat{x}$ 共轭的动量算符，满足 $[\hat{q}, \hat{x}] = \mathrm{i}$。流算符由 $\hat{I} \equiv \hat{x}/L$ 给出。

我们也可以把 $q$ 视为坐标，而把 $-x$ 视为相应的“动量”。在这种观点下，哈密顿量 (2.4.14) 给出如下拉格朗日量，它是式 (2.4.13) 的对偶形式（dual form）：
\begin{equation*}
L_{\mathrm{CL}}^d = \frac{L}{2}\dot{q}^2 - \frac{1}{2C}q^2 \tag{2.4.15}
\end{equation*}
如图 2.10 所示，在加入一个电阻器之后，经典运动方程变为
\begin{equation*}
V = \frac{q}{C} = L\frac{\mathrm{d}I}{\mathrm{d}t}, \qquad\qquad I = -\frac{\mathrm{d}q}{\mathrm{d}t} - V/R
\end{equation*}
由此得到
\begin{equation*}
C\frac{\mathrm{d}^2x}{\mathrm{d}t^2} = -\frac{x}{L} - \frac{1}{R}\frac{\mathrm{d}x}{\mathrm{d}t} \tag{2.4.16}
\end{equation*}
这正是摩擦系数为 $\gamma = 1/R$ 的粒子的运动方程。这样的系统可以用一个双时作用量描述（见式 (2.4.12)）：
\begin{equation*}
S_{\mathrm{RCL}} = \int \mathrm{d}t\,\left(\frac{1}{2}C\dot{x}^2 - \frac{1}{2L}x^2\right) + \mathrm{i}\int \mathrm{d}t\,\mathrm{d}t'\;\frac{(x(t) - x(t'))^2}{4\pi R(t - t')^2} \tag{2.4.17}
\end{equation*}
RCL 电路的许多量子动力学性质都可以由上述作用量计算出来。

如果我们计入电荷量子化，那么电荷为 $q = e^{*} \times \text{整数}$。由于电荷 $q$ 也是 $x$ 的动量，动量的量子化意味着粒子生活在一个圆上。电荷（即动量）量子 $e^{*}$ 意味着 $x$ 与 $x + \frac{2\pi}{e^{*}}$ 应当被视为同一个点。

由于电荷量子化蕴含 $x$ 的周期性，描述带电荷量子化的 RCL 电路的作用量必须是 $x$ 的周期函数。

有许多办法可以使式 (2.4.17) 周期化。最简单的周期双时作用量之一由下式给出：
\begin{equation*}
S_{\mathrm{RCL}} = \int \mathrm{d}t\,\left(\frac{C}{2}\dot{x}^2 - \frac{1 - \cos(e^{*}x)}{Le^{*2}}\right) + \mathrm{i}\int \mathrm{d}t\,\mathrm{d}t'\;\frac{\sin^2\!\left(\frac{e^{*}}{2}[x(t) - x(t')]\right)}{\pi R e^{*2}(t - t')^2} \tag{2.4.18}
\end{equation*}
我们注意到，在 $e^{*} \to 0$ 极限下式 (2.4.18) 变为式 (2.4.17)。式 (2.4.18) 描述了带电荷量子化的 RCL 电路的量子动力学性质。

\subsection{耗散与涨落之间的关系}

\begin{keypoints}
\item 耗散与涨落总是同时出现。我们可以由其一确定另一个。
\end{keypoints}

双时作用量 (2.4.12) 与 (2.4.17) 适合于计算平衡性质。让我们计算 $x$ 与速度 $v = \dot{x}$ 的平衡量子涨落。$x$ 的涨落由噪声功率谱（noise power spectrum）表征。为了定义噪声功率谱，让我们先考虑一个经典涨落量 $x(t)$。总功率（包括所有频率）定义为
\begin{equation*}
P_{\mathrm{tot}} \equiv \int_0^{t_\infty} \mathrm{d}t\,|x(t)|^2/t_\infty = 2\int_0^{\infty} \frac{\mathrm{d}\omega}{2\pi}\,x_{-\omega}x_{\omega}/t_\infty
\end{equation*}
其中 $x_\omega = \int_0^{t_\infty} \mathrm{d}t\;x(t)\,\mathrm{e}^{\mathrm{i}\omega t}$，而 $t_\infty \to \infty$。于是，我们可以把 $x_{-\omega}x_\omega/t_\infty$ 认作噪声功率谱。量子噪声功率谱则由对 $x_{-\omega}x_\omega/t_\infty$ 取量子平均而得到：\footnote{这里我们用了
\begin{align*}
2\langle x_{-\omega}x_\omega\rangle/t_\infty &= \int_0^{t_\infty} \mathrm{d}t_1\,\mathrm{d}t_2\;\left[\langle x(t_1)x(t_2)\rangle + \langle x(t_2)x(t_1)\rangle\right]\mathrm{e}^{\mathrm{i}\omega(t_1 - t_2)}/t_\infty\\
&= \int_0^{t_\infty} \mathrm{d}t_1\,\mathrm{d}t_2\;\left[\langle T(x(t_1)x(t_2))\rangle + \langle T(x(t_1)x(t_2))\rangle^{*}\right]\mathrm{e}^{\mathrm{i}\omega(t_1 - t_2)}/t_\infty\\
&= \int_0^{t_\infty} \mathrm{d}t_1\,\mathrm{d}t_2\;\left[\langle T(x(t_1)x(t_2))\rangle + \langle T(x(t_2)x(t_1))\rangle^{*}\right]\mathrm{e}^{\mathrm{i}\omega(t_1 - t_2)}/t_\infty\\
&= 2\,\mathrm{Re}\int_{-\infty}^{+\infty} \mathrm{d}t\;\langle T(x(t)x(0))\rangle\,\mathrm{e}^{\mathrm{i}\omega t} = -2\mathrm{Im}G_\omega
\end{align*}}
\begin{equation*}
P(\omega) = 2\langle x_{-\omega}x_\omega\rangle/t_\infty = -2\mathrm{Im}G_\omega \tag{2.4.19}
\end{equation*}
由 $\mathrm{Im}D$ 与 $\mathrm{Im}G$ 之间的如下关系（见式 (2.2.42)）：
\begin{equation*}
\mathrm{Im}G(\omega) = \frac{\mathrm{Im}D}{\tanh(\beta\omega/2)}, \tag{2.4.20}
\end{equation*}
我们看到 $\mathrm{Im}G$（从而 $P(\omega)$）可以由响应函数 $D$ 确定。

响应函数可以由运动方程算出。对于带摩擦的粒子，运动方程为
\begin{equation*}
m\ddot{x} = -Kx - \gamma\dot{x} + f(t) \tag{2.4.21}
\end{equation*}
其中 $f$ 是外力。运动方程的解可以形式地写为
\begin{equation*}
x(t) = -\left(\frac{1}{-m\frac{\mathrm{d}^2}{\mathrm{d}t^2} - \gamma\frac{\mathrm{d}}{\mathrm{d}t} - K}\right)f
\end{equation*}
这个解可以解释为 $x$ 对微扰 $\delta H = -f(t)x$ 的响应：$x(t) = \int \mathrm{d}t'\;D(t - t')(-f(t')) \equiv -Df$（见式 (2.2.3)，取 $O_1 = O_2 = x$）。因此，运动方程决定了 $x$ 的如下响应函数：
\begin{equation*}
D = \frac{1}{-m\frac{\mathrm{d}^2}{\mathrm{d}t^2} - \gamma\frac{\mathrm{d}}{\mathrm{d}t} - K}
\end{equation*}
现在清楚了：响应函数的 $\mathrm{Im}D$ 描写运动粒子的摩擦（即耗散），而关联函数的 $\mathrm{Im}G$ 描写涨落的功率谱。式 (2.4.20) 是耗散与涨落之间的一个直接关系。

由 $D$ 的表达式，我们求得 $x$ 涨落的功率谱为（见式 (2.4.19)）
\begin{align*}
P(\omega) &= -2\left(\tanh(\beta\omega/2)\right)^{-1}\mathrm{Im}D\\
&= \frac{1}{\tanh(\omega\hbar/2T)}\;\frac{2\gamma\omega\hbar}{(m\omega^2 - K)^2 + \gamma^2\omega^2} \tag{2.4.22}
\end{align*}
速度涨落的功率谱为
\begin{equation*}
P^v(\omega) = \omega^2 P(\omega) = \frac{1}{\tanh(\omega\hbar/2T)}\;\frac{2\gamma\omega\hbar}{(m\omega - \frac{K}{\omega})^2 + \gamma^2}
\end{equation*}
我们注意到，运动方程 (2.4.21) 与式 (2.4.16) 给出的 RCL 电路的运动方程完全相同，只需把 $(C, L^{-1}, R^{-1})$ 换成 $(m, K, \gamma)$。电压由 $V = \dot{x}$ 给出。于是，RCL 电路中电压涨落的功率谱为
\begin{equation*}
P^V(\omega) = \frac{1}{\tanh(\hbar\beta\omega/2)}\;\frac{2\hbar R\omega}{(RC\omega - \frac{R}{L\omega})^2 + 1} \tag{2.4.23}
\end{equation*}
当 $\hbar\omega \ll T$ 时，我们有
\begin{equation*}
P_{\mathrm{c}}^V(\omega) = \frac{4TR}{(RC\omega - \frac{R}{L\omega})^2 + 1}
\end{equation*}
这就是 RCL 系统的经典噪声谱。当 $L = \infty$ 且 $C = 0$ 时，上式变为一个电阻器两端的噪声谱，即 $P_{\mathrm{c}}^V(\omega) = 4TR$。

有趣的是，有限温度噪声谱 $P^V$ 与零温噪声谱 $P_0^V$ 之间有如下紧密的联系：
\begin{equation*}
P^V(\omega) = \frac{1}{\tanh(\hbar|\omega|/2T)}\,P_0^V(\omega) \tag{2.4.24}
\end{equation*}
$P^V(\omega)$ 也可以由经典噪声谱 $P_{\mathrm{c}}^V(\omega)$ 得到：
\begin{equation*}
P^V(\omega) = \frac{\hbar\omega}{2T\tanh(\hbar\omega/2T)}\,P_{\mathrm{c}}^V(\omega)
\end{equation*}
上述关系非常普遍，适用于所有谐系统（见下面的问题 2.4.5）。

\subsection*{问题 2.4.4}

\textbf{布朗运动（Brownian motion）：}
\begin{enumerate}
\item[(a)] 一个质量为 $m$ 的粒子受到由 $\gamma$ 描述的摩擦。如果在 $x = 0$ 处释放粒子，经过时间 $t$ 之后它能够游荡到多远？（即在时刻 $t$ 计算 $\bar{x} = \sqrt{\langle x^2\rangle}$。）（提示：考虑 $\omega = 0$ 处速度的功率谱。）
\item[(b)] 对悬浮在水中的一个塑料球，求出（或估计）1 分钟之后 $\bar{x}$ 的数值。球的直径为 1 cm。再对水中的一粒花粉重复这一计算，我们把花粉建模为用同样塑料做成的、尺寸小 10000 倍的球。20°C 时水的黏度为 0.01 泊（1 泊 = 1 达因·秒/cm$^2$）。
\end{enumerate}

\subsection*{问题 2.4.5}

直接考虑一个具有任意 $g_i$ 与 $\Omega_i$ 的耦合振子系统（见式 (2.4.8)），证明式 (2.4.24)。（提示：考虑振子在 $T = 0$ 与 $T \neq 0$ 时的 $\langle x^2\rangle$。）

\subsection{随机微分方程的路径积分描述}

\begin{keypoints}
\item 随机微分方程之解的统计性质可以用路径积分来描述。
\end{keypoints}

按照运动方程 (2.4.21)，如果把粒子在某处释放，那么当 $f = 0$ 时，粒子经过很长时间后最终会停在 $x = 0$ 处。

然而，耗散/涨落定理告诉我们，上述图像并不正确，粒子并非只停留在 $x = 0$ 处。摩擦的存在蕴含着涨落的存在，因此粒子会围绕 $x = 0$ 涨落。要得到一个描写耗散系统的更正确的运动方程，我们需要加入一个产生涨落的项。一种办法是让力项 $f(t)$ 随时间随机涨落。现在有两个问题。第一，随机力项能模拟涨落吗——特别是量子涨落？如果答案是可以，那么第二，我们如何找到能产生正确涨落的随机力的概率分布？为解决这个问题，我们需要计算式 (2.4.21) 的随机解的平均关联 $\langle x(t)x(0)\rangle$，然后考察是否可能选出一个恰当的 $f(t)$ 分布，使其重现功率谱 (2.4.22)。

我们注意到
\begin{equation*}
\langle x(t_b)x(t_a)\rangle = \int \mathcal{D}[f(t)]\,\mathcal{D}[x(t)]\;x(t_b)\,x(t_a)\,P[f(t)]\,\delta[x(t) - x^f(t)]
\end{equation*}
这里 $x^f = \mathcal{K}^{-1}f$ 是式 (2.4.21) 的解，其中 $\mathcal{K} = m\frac{\mathrm{d}^2}{\mathrm{d}t^2} + \gamma\frac{\mathrm{d}}{\mathrm{d}t} + K$，而 $P[f(t)]$ 是 $f(t)$ 的概率分布（probability distribution）。我们注意到 $\mathcal{K} = -D^{-1}$。由于算符 $\mathcal{K}$ 不依赖于 $f$，我们有
\begin{equation*}
\delta[x(t) - x^f(t)] \propto \delta[\mathcal{K}x(t) - f(t)]
\end{equation*}
于是，
\begin{align*}
\langle x(t_b)x(t_a)\rangle &= \frac{\int \mathcal{D}[f(t)]\,\mathcal{D}[x(t)]\;x(t_b)x(t_a)\,P[f(t)]\,\delta[\mathcal{K}x(t) - f(t)]}{\int \mathcal{D}[f(t)]\,\mathcal{D}[x(t)]\;P[f(t)]\,\delta[\mathcal{K}x(t) - f(t)]}\\
&= \frac{\int \mathcal{D}[f(t)]\,\mathcal{D}[x(t)]\,\mathcal{D}[\lambda(t)]\;x(t_b)x(t_a)\,P[f(t)]\,\mathrm{e}^{\mathrm{i}\int \lambda[\mathcal{K}x - f]}}{\int \mathcal{D}[f(t)]\,\mathcal{D}[x(t)]\,\mathcal{D}[\lambda(t)]\;P[f(t)]\,\mathrm{e}^{\mathrm{i}\int \lambda[\mathcal{K}x - f]}}
\end{align*}
这是一个相当有趣的结果。我们已经看到，路径积分可以描述哈密顿量——量子系统中的一个线性算符；它也能描述经典及量子统计系统中的热平均。从上面我们看到，路径积分甚至可以描述一个随机微分方程（random differential equation）。

现在让我们假设随机项的分布是高斯型的：
\begin{equation*}
P[f(t)] \propto \mathrm{e}^{-\frac{1}{2}\int f\mathcal{V}f}
\end{equation*}
其中 $\int f\mathcal{V}f \equiv \int \mathrm{d}t\,\mathrm{d}t'\;f(t)\mathcal{V}(t,t')f(t')$。我们可以把 $f$ 与 $\lambda$ 依次积掉，如下：
\begin{align*}
\langle x(t_b)x(t_a)\rangle &= \frac{\int \mathcal{D}[x(t)]\,\mathcal{D}[\lambda(t)]\;x(t_b)x(t_a)\,\mathrm{e}^{\int \mathrm{i}\lambda\mathcal{K}x - \frac{1}{2}\lambda\mathcal{V}^{-1}\lambda}}{\int \mathcal{D}[x(t)]\,\mathcal{D}[\lambda(t)]\;\mathrm{e}^{\int \mathrm{i}\lambda\mathcal{K}x - \frac{1}{2}\lambda\mathcal{V}^{-1}\lambda}}\\
&= \frac{\int \mathcal{D}[x(t)]\;x(t_b)x(t_a)\,\mathrm{e}^{-\int \frac{1}{2}x\mathcal{K}^{\top}\mathcal{V}\mathcal{K}x}}{\int \mathcal{D}[x(t)]\;\mathrm{e}^{-\int \frac{1}{2}x\mathcal{K}^{\top}\mathcal{V}\mathcal{K}x}}\\
&= (\mathcal{K}^{\top}\mathcal{V}\mathcal{K})^{-1}(t_b - t_a)
\end{align*}
这里，算符 $\mathcal{K}^{\top}\mathcal{V}\mathcal{K}(t)$ 的逆被理解为 $\mathcal{G}(t - t')$，使得 $\mathcal{K}^{\top}\mathcal{V}\mathcal{K}(t)\mathcal{G}(t - t') = \delta(t - t')$。我们注意到 $\langle x(t_b)x(t_a)\rangle = \langle x(t_a)x(t_b)\rangle$，所以它的傅里叶变换是实的，并且等于 $x$ 涨落功率谱的一半（见式 (2.4.19)）。于是，为了重现 $x$ 的涨落，我们要求
\begin{equation*}
\mathcal{K}^{\top}\mathcal{V}\mathcal{K} = \frac{\tanh(\omega\hbar/2T)\left[(-m\omega^2 + K)^2 + \omega^2\gamma^2\right]}{\gamma\omega\hbar}
\end{equation*}
或者（在频率空间中）
\begin{equation*}
\mathcal{V} = \frac{\tanh(\omega\hbar/2T)}{\gamma\omega\hbar}
\end{equation*}
我们看到，即便是量子噪声也可以用随机的经典方程来模拟。特别地，我们注意到随机力项只依赖于温度 $T$ 与摩擦系数 $\gamma$。更高的温度或更大的摩擦会导致更强的随机力。经典噪声（在 $\hbar \to 0$ 极限下）可以用随机力的一个简单概率分布模拟如下：
\begin{equation*}
P[f(t)] \propto \mathrm{e}^{-\frac{1}{2}\int \mathrm{d}t\;f\mathcal{V}f}, \qquad \mathcal{V} = \frac{1}{2\gamma T}
\end{equation*}
它在时间上是无关联的（即 $f(t)$ 与 $f(t')$ 独立地涨落）。

我们想指出，随机微分方程理论是一个非常庞大的领域，有着广泛而重要的应用。随机微分方程被用于描述材料生长、化学与生物过程，甚至经济系统。随机微分方程的路径积分表示使我们能够把为路径积分发展出的图像与概念——如对称性破缺、相变、重整化群等——应用到随机微分方程上。

% ==================================================================
% 新增术语（ch2_e 新增，供术语表追加）：
%   heat bath 热库；friction coefficient 摩擦系数；effective mass 有效质量；
%   double-time action / double-time path integral 双时作用量 / 双时路径积分；
%   closed-time path integral 闭合时间路径积分；Schwinger–Keldysh formalism Schwinger–Keldysh 形式；
%   time reversal 时间反演；equilibrium properties 平衡性质；non-equilibrium 非平衡；
%   capacitor / capacitance 电容器 / 电容；inductance 电感；resistor 电阻器；
%   RCL / CL circuit RCL / CL 电路；charge quantization 电荷量子化；
%   charge operator 电荷算符；current operator 流算符；dual form 对偶形式；
%   decaying process 衰减过程；decay (of a state) 衰变；
%   noise power spectrum 噪声功率谱；power spectrum 功率谱；noise 噪声；
%   Brownian motion 布朗运动；viscosity 黏度；poise 泊；dyne 达因；
%   random differential equation 随机微分方程；random force 随机力；
%   probability distribution 概率分布；discontinuity 不连续性
% ==================================================================
